\begin{lexlean}{Divisibility}
\useglossary{lexlean.std.nat@1.0.0}
\useglossary{peano.arith@1.0.0}
\importmodule{Addition}
\title{Natural number divisibility}

\begin{section}{definition}
\heading{Divisibility}
  \begin{predicatedefinition}{divides}{peano.arith::divides}
  \noaxioms
  For every natural number \(n\); natural number \(m\), \(n\) divides \(m\) holds exactly when there exists a natural number \(k\) such that \(m = n * k\).
  \end{predicatedefinition}
\end{section}

\begin{section}{properties}
\heading{Divisibility : order}
  \begin{theorem}{divides-refl}
  \noaxioms
  For every natural number \(n\), \(n\) divides \(n\).
  \begin{proof}
  Use \(1\) as the witness.
  Close the goal with \(eqsymm(mulone(n))\).
  \end{proof}
  \end{theorem}

  \begin{theorem}{one-divides}
  \noaxioms
  For every natural number \(n\), \(1\) divides \(n\).
  \begin{proof}
  Use \(n\) as the witness.
  Close the goal with \(eqsymm(onemul(n))\).
  \end{proof}
  \end{theorem}

  \begin{theorem}{divides-trans}
  \exactaxioms{propext}
  For every natural number \(a\) and natural number \(b\) and natural number \(c\), if \(a\) divides \(b\) and \(b\) divides \(c\), then \(a\) divides \(c\).
  \begin{proof}
  Assume \(h\).
  \begin{cases}{h}
  \begin{case}{lexlean.core::and-intro}
  \bind{hab;hbc}
  \begin{cases}{hab}
  \begin{case}{lexlean.core::exists-intro}
  \bind{k;hk}
  \begin{cases}{hbc}
  \begin{case}{lexlean.core::exists-intro}
  \bind{l;hl}
  Use \(k * l\) as the witness.
  \begin{rewrite}{goal}
  \forward{hl}
  \forward{hk}
  \end{rewrite}
  Close the goal with \(mulassoc(a, k, l)\).
  \end{case}
  \end{cases}
  \end{case}
  \end{cases}
  \end{case}
  \end{cases}
  \end{proof}
  \end{theorem}
\end{section}
\end{lexlean}
